% Generated draft: requires mathematical and source review.
\documentclass[11pt]{article}
\usepackage[margin=25mm]{geometry}
\usepackage{fontspec}
\setmainfont{Latin Modern Roman}
\usepackage{amsmath,amssymb,mathtools,unicode-math}
\setmathfont{Latin Modern Math}
\usepackage{enumitem,needspace,xurl,hyperref}
\setlength{\parindent}{0pt}
\setlength{\parskip}{5pt}
\setlength{\emergencystretch}{3em}
\newcommand{\N}{\mathbb{N}}
\newcommand{\Z}{\mathbb{Z}}
\newcommand{\Q}{\mathbb{Q}}
\newcommand{\R}{\mathbb{R}}
\newcommand{\C}{\mathbb{C}}
\newcommand{\F}{\mathbb{F}}
\newcommand{\A}{\mathbb{A}}
\newcommand{\PP}{\mathbb{P}}
\newcommand{\E}{\mathbb{E}}
\newcommand{\eps}{\varepsilon}
\newcommand{\dd}{\,\mathrm{d}}
\begin{document}
\section*{UnsolvedMath ID: 30000495}
\textbf{Source problem code:} catalogue entry
\subsection{Definitions and mathematical statement}
Let \(k\) be a finite field of characteristic \(p\), and let \(G\) be a compact \(p\)-adic analytic group with no elements of order \(p\). Put \(d=\dim G\). In the usual Iwasawa-algebra convention, the notation \(kG\) here means the completed group algebra
\[
A=k[[G]]=\varprojlim_U k[G/U],
\]
where \(U\) ranges over open normal subgroups. This convention must be checked; the ordinary group ring would be a different object. The standard canonical dimension of a nonzero finitely generated module is \(d-j_A(M)\), where \(j_A(M)=\inf\{i:\operatorname{Ext}_A^i(M,A)\ne0\}\).

Let \(\mathcal F_t\) be the source's exact category of finitely generated modules of canonical dimension at most \(t\), for \(0\leq t\leq d\). Inclusion gives
\[
\theta_t:K_0(\mathcal F_t)\longrightarrow K_0(\mathcal F_d).
\]
Here \(K_0\) uses short-exact-sequence relations for these module categories, and
\(\operatorname{rank}(\operatorname{im}\theta_t)=\dim_\Q(\operatorname{im}\theta_t\otimes_\Z\Q)\).

The first question asks to improve the reported upper bound by the number of Galois orbits of conjugacy classes of finite-order elements \(g\in G\) with \(\dim C_G(g)\leq t\). In symbols, if \(\mathcal C_t\) denotes those conjugacy classes and \(\mathscr G_k\) the intended Galois action, the bound being refined is read as
\[
\operatorname{rank}(\operatorname{im}\theta_t)\leq|\mathscr G_k\backslash\mathcal C_t|.
\]
Since the torsion orders are prime to \(p\), the standard finite-field action is generated by powering \(g\mapsto g^{|k|}\); its exact convention requires source confirmation. The source does not specify a sharper proposed formula or what improvement is sought.

The second question asks for an analogue when \(k\) is replaced by \(\Z_p\). In the completed convention the algebra becomes \(\Z_p[[G]]\); its homological dimension convention and canonical-dimension filtration need to be stated afresh, rather than copying \(d\) and the finite-field orbit count unchanged. The record gives no integer-coefficient bound.

The supplied assessment checked 2026-08-21 says both questions remain open. Its additional reference is labeled as a mixed \(K(\pi,1)\) paper, whose relevance to these module Grothendieck-group rank questions is not explained; it must not be treated as a matched theorem proving or disproving either requested refinement.
\subsection{Short English statement}
Sharpen the stated Galois-orbit upper bound for the images of canonical-dimension filtration maps on Grothendieck groups, and formulate the corresponding result over \(\Z_p\). The completed-ring convention and the integer-coefficient filtration must be checked.
\subsection{Sources}
\begin{itemize}
\item Local UnsolvedMath registry, permanent ID 30000495, catalogue entry, \emph{Rank Bounds for Modular-Representation Grothendieck Groups}.
\item \emph{Pro-\(p\) Extensions of Global Fields and pro-\(p\) Groups} (2007), Oberwolfach Reports, pp. 1463--1536, DOI \url{https://doi.org/10.4171/owr/2006/25}.
\item Supplied reference ``Schmidt on the mixed K-pi-1 case'': \url{https://arxiv.org/abs/0801.2103}. Its specific relevance to the extracted rank-bound questions is not supplied.
\item Supplied direct statement-verification claim restores the hypotheses, filtration map, bound, and coefficient extension; supplied literature triage checked 2026-08-21.
\end{itemize}
\end{document}
